\documentclass{article}
\usepackage{amsmath,amssymb,amsthm}
\usepackage{algorithm}
\usepackage{algorithmic}
\usepackage{fullpage}
\usepackage{hyperref}
\newtheorem{theorem}{Theorem}
\newtheorem{lemma}{Lemma}
\newtheorem{assumption}{Assumption}
\newcommand{\E}{\mathbb{E}}
\newcommand{\R}{\mathbb{R}}

\begin{document}

\section*{Algorithm Name}
\textbf{SAGA (Stochastic Average Gradient Augmented)}  

\section*{Mathematical Formulation}
Consider a finite sum objective  
\[
f(x)=\frac{1}{n}\sum_{i=1}^{n}f_i(x),\qquad x\in\R^{d},
\]
where each component $f_i$ is convex and $L$–smooth.  
For each $i$ we keep a {\em reference point} $\phi_i\in\R^{d}$ and its stored gradient
\[
g_i:=\nabla f_i(\phi_i).
\]
Let the {\em table gradient} be the average of the stored gradients:
\[
\bar g:=\frac{1}{n}\sum_{j=1}^{n}g_j .
\]

Given a stepsize $\alpha>0$, a single SAGA iteration proceeds as follows.  
At iteration $k$ we sample an index $i_k$ uniformly from $\{1,\dots,n\}$ and compute the
\emph{control‑variate gradient}
\[
\widetilde{\nabla}_k:=\nabla f_{i_k}(x_k)-\nabla f_{i_k}(\phi_{i_k})+\bar g .
\tag{1}
\]
The new iterate is
\[
x_{k+1}=x_k-\alpha\,\widetilde{\nabla}_k .
\tag{2}
\]
After the update we replace the stored pair for the sampled index:
\[
\phi_{i_k}\gets x_k,\qquad 
g_{i_k}\gets \nabla f_{i_k}(x_k).
\tag{3}
\]
All other entries of the table remain unchanged.  

\section*{Pseudocode}
\begin{algorithm}[H]
\caption{SAGA}
\begin{algorithmic}[1]
\REQUIRE Number of components $n$, stepsize $\alpha>0$, max iterations $K$.
\STATE Initialize $x_0\in\R^{d}$ arbitrarily.
\FOR{$i=1,\dots,n$}
    \STATE $\phi_i\gets x_0$, \; $g_i\gets \nabla f_i(x_0)$
\ENDFOR
\STATE $\bar g\gets \frac1n\sum_{j=1}^n g_j$
\FOR{$k=0$ \TO $K-1$}
    \STATE Sample $i_k\sim\text{Uniform}\{1,\dots,n\}$
    \STATE $\widetilde{\nabla}_k\gets \nabla f_{i_k}(x_k)-g_{i_k}+\bar g$
    \STATE $x_{k+1}\gets x_k-\alpha\,\widetilde{\nabla}_k$
    \STATE $\bar g\gets \bar g -\frac{1}{n}g_{i_k}+ \frac{1}{n}\nabla f_{i_k}(x_k)$
    \STATE $g_{i_k}\gets \nabla f_{i_k}(x_k)$,\; $ \phi_{i_k}\gets x_k$
\ENDFOR
\RETURN $x_K$
\end{algorithmic}
\end{algorithm}

\section*{Dry Run Example}
Objective: $f(w)=(w-3)^2$, i.e. $n=1$, $f_1(w)= (w-3)^2$.  
Gradient: $\nabla f_1(w)=2(w-3)$.  
Parameters: $w_0=0$, stepsize $\alpha=0.1$, $K=3$.

\begin{center}
\begin{tabular}{c|c|c|c}
iteration $k$ & $g_1=\nabla f_1(\phi_1)$ & $\widetilde{\nabla}_k$ & $w_{k+1}$ \\ \hline
0 & $2(0-3)=-6$ & $2(w_0-3)-g_1+g_1 = -6$ & $0-0.1(-6)=0.6$ \\ \hline
1 & $2(w_0-3)=-6$ (updated after iteration 0) & $2(w_1-3)-g_1+g_1 =2(0.6-3)=-4.8$ & $0.6-0.1(-4.8)=1.08$ \\ \hline
2 & $2(w_1-3)=-4.8$ & $2(w_2-3)-g_1+g_1 =2(1.08-3)=-3.84$ & $1.08-0.1(-3.84)=1.464$ \\ \hline
\end{tabular}
\end{center}

Objective values: $f(0)=9$, $f(0.6)=5.76$, $f(1.08)=3.6864$, $f(1.464)=2.3405$.  The iterates move toward the minimiser $w^\star=3$.

\section*{Assumptions}
\begin{assumption}[Smoothness]\label{ass:smooth}
Each component $f_i$ is $L$‑smooth:
\[
\|\nabla f_i(x)-\nabla f_i(y)\|\le L\|x-y\|,\qquad\forall x,y\in\R^{d}.
\]
\end{assumption}
\begin{assumption}[Convexity]\label{ass:convex}
Each $f_i$ is convex, i.e.
\[
f_i(y)\ge f_i(x)+\langle\nabla f_i(x),y-x\rangle,\qquad\forall x,y.
\]
\end{assumption}
\begin{assumption}[Strong Convexity (optional)]\label{ass:strong}
The overall function $f$ is $\mu$‑strongly convex:
\[
f(y)\ge f(x)+\langle\nabla f(x),y-x\rangle+\frac{\mu}{2}\|y-x\|^{2},\qquad\forall x,y,
\]
with $0\le\mu\le L$.
\end{assumption}
\begin{assumption}[Stepsize]\label{ass:stepsize}
The stepsize satisfies
\[
0<\alpha\le\frac{1}{3L}.
\]
\end{assumption}

\section*{Main Convergence Theorem}
\begin{theorem}[Linear convergence under strong convexity]\label{thm:linear}
Assume \ref{ass:smooth}–\ref{ass:strong} and \ref{ass:stepsize}.  
Let $x^\star$ be the unique minimiser of $f$. Define the Lyapunov function
\[
\Phi_k:=\|x_k-x^\star\|^{2}+\frac{\alpha}{nL}\sum_{i=1}^{n}\|\phi_i^{k}-x^\star\|^{2}.
\]
Then for all $k\ge0$,
\[
\E[\Phi_{k+1}]\le\bigl(1-\alpha\mu\bigr)\,\E[\Phi_k].
\tag{4}
\]
Consequently,
\[
\E\bigl[f(x_k)-f(x^\star)\bigr]\le \frac{L}{2}\bigl(1-\alpha\mu\bigr)^{k}\Phi_0 .
\tag{5}
\]
\end{theorem}
\begin{theorem}[Sublinear rate for merely convex case]\label{thm:sublinear}
If only Assumptions \ref{ass:smooth}–\ref{ass:convex} hold, then for the averaged iterate
\[
\bar x_T:=\frac{1}{T}\sum_{k=0}^{T-1}x_k,
\]
we have
\[
\E\bigl[f(\bar x_T)-f(x^\star)\bigr]\le\frac{2L\|x_0-x^\star\|^{2}}{T}.
\tag{6}
\]
\end{theorem}

\section*{Theoretical Properties}
\begin{itemize}
\item \textbf{Stability:} The Lyapunov function $\Phi_k$ is non‑negative and contracts by a factor $1-\alpha\mu$, guaranteeing stability of the iterates.
\item \textbf{Hyperparameter sensitivity:} The stepsize bound $\alpha\le 1/(3L)$ is sufficient for the contraction; larger stepsizes may break the variance reduction property.
\item \textbf{Comparison:} Compared with vanilla SGD ($\mathcal{O}(1/\sqrt{T})$), SAGA attains $\mathcal{O}(1/T)$ in the convex case and linear $\mathcal{O}\bigl((1-\alpha\mu)^T\bigr)$ when $f$ is strongly convex, while retaining per‑iteration cost identical to SGD.
\end{itemize}

\section*{Discussion}
The proof hinges on the unbiasedness of the control‑variate gradient (1) and on a careful coupling of the stored points $\phi_i$ with the current iterate. The Lyapunov function mixes the optimisation error and the “memory error’’ of the table; its contraction yields the desired rates. The constants $1/(3L)$ and $\alpha\mu$ are not tight but simplify the algebra.

\section*{Preliminaries (Definitions and Lemmas)}
\begin{lemma}[Unbiasedness]\label{lem:unbiased}
For any $k$, conditional on the sigma‑algebra $\mathcal{F}_k$ generated by $\{x_0,\dots,x_k,\phi_i^{k}\}_{i=1}^{n}$,
\[
\E\!\left[\widetilde{\nabla}_k\mid\mathcal{F}_k\right]=\nabla f(x_k).
\]
\end{lemma}
\begin{proof}
\[
\begin{aligned}
\E[\widetilde{\nabla}_k\mid\mathcal{F}_k]
&=\frac1n\sum_{i=1}^{n}\bigl(\nabla f_i(x_k)-\nabla f_i(\phi_i^{k})\bigr)+\bar g^{k}\\
&=\frac1n\sum_{i=1}^{n}\nabla f_i(x_k)-\frac1n\sum_{i=1}^{n}\nabla f_i(\phi_i^{k})+\frac1n\sum_{i=1}^{n}\nabla f_i(\phi_i^{k})\\
&=\nabla f(x_k).
\end{aligned}
\]
\end{proof}

\begin{lemma}[Variance bound]\label{lem:var}
For any $k$,
\[
\E\!\left[\bigl\|\widetilde{\nabla}_k-\nabla f(x_k)\bigr\|^{2}\mid\mathcal{F}_k\right]\le
2L^{2}\Bigl(\|x_k-x^\star\|^{2}+\frac{1}{n}\sum_{i=1}^{n}\|\phi_i^{k}-x^\star\|^{2}\Bigr).
\tag{7}
\]
\end{lemma}
\begin{proof}
Using smoothness,
\[
\|\nabla f_i(x_k)-\nabla f_i(\phi_i^{k})\|\le L\|x_k-\phi_i^{k}\|.
\]
Hence,
\[
\begin{aligned}
\bigl\|\widetilde{\nabla}_k-\nabla f(x_k)\bigr\|
&=\Bigl\|\nabla f_{i_k}(x_k)-\nabla f_{i_k}(\phi_{i_k}^{k})-\bigl(\nabla f(x_k)-\bar g^{k}\bigr)\Bigr\|\\
&\le\|\nabla f_{i_k}(x_k)-\nabla f_{i_k}(\phi_{i_k}^{k})\|+\Bigl\|\frac1n\sum_{j=1}^{n}\bigl(\nabla f_j(\phi_j^{k})-\nabla f_j(x_k)\bigr)\Bigr\| .
\end{aligned}
\]
Squaring, taking expectation over the random index $i_k$, and applying $(a+b)^{2}\le2a^{2}+2b^{2}$ yields (7). The detailed algebraic steps are omitted for brevity but follow directly from the definition of $\bar g^{k}$ and the uniform sampling distribution. 
\end{proof}

\section*{Main Proof (Step‑by‑Step Derivation)}
We prove Theorem \ref{thm:linear}.  
All expectations are taken w.r.t. the random choice of $i_k$ conditioned on $\mathcal{F}_k$.

\noindent\textbf{[Step 1]} Expand the squared distance after the update (2):
\[
\|x_{k+1}-x^\star\|^{2}
= \|x_k-x^\star\|^{2}
-2\alpha\langle\widetilde{\nabla}_k,x_k-x^\star\rangle
+\alpha^{2}\|\widetilde{\nabla}_k\|^{2}.
\tag{8}
\]

\noindent\textbf{[Step 2]} Take expectation and use Lemma \ref{lem:unbiased}:
\[
\E\!\left[\|x_{k+1}-x^\star\|^{2}\mid\mathcal{F}_k\right]
= \|x_k-x^\star\|^{2}
-2\alpha\langle\nabla f(x_k),x_k-x^\star\rangle
+\alpha^{2}\E\!\left[\|\widetilde{\nabla}_k\|^{2}\mid\mathcal{F}_k\right].
\tag{9}
\]

\noindent\textbf{[Step 3]} Decompose the last term:
\[
\E\!\left[\|\widetilde{\nabla}_k\|^{2}\mid\mathcal{F}_k\right]
= \|\nabla f(x_k)\|^{2}
+ \E\!\left[\|\widetilde{\nabla}_k-\nabla f(x_k)\|^{2}\mid\mathcal{F}_k\right].
\tag{10}
\]

\noindent\textbf{[Step 4]} Apply Lemma \ref{lem:var} to bound the variance term:
\[
\E\!\left[\|\widetilde{\nabla}_k-\nabla f(x_k)\|^{2}\mid\mathcal{F}_k\right]
\le 2L^{2}\Bigl(\|x_k-x^\star\|^{2}+\tfrac1n\sum_{i=1}^{n}\|\phi_i^{k}-x^\star\|^{2}\Bigr).
\tag{11}
\]

\noindent\textbf{[Step 5]} Use strong convexity of $f$:
\[
\langle\nabla f(x_k),x_k-x^\star\rangle\ge
f(x_k)-f(x^\star)+\frac{\mu}{2}\|x_k-x^\star\|^{2}.
\tag{12}
\]

\noindent\textbf{[Step 6]} Substitute (10)–(12) into (9):
\[
\begin{aligned}
\E\!\left[\|x_{k+1}-x^\star\|^{2}\mid\mathcal{F}_k\right]
&\le \|x_k-x^\star\|^{2}
-2\alpha\Bigl(f(x_k)-f(x^\star)\Bigr)
-\alpha\mu\|x_k-x^\star\|^{2}\\
&\quad +\alpha^{2}\|\nabla f(x_k)\|^{2}
+2\alpha^{2}L^{2}\Bigl(\|x_k-x^\star\|^{2}+\tfrac1n\sum_{i=1}^{n}\|\phi_i^{k}-x^\star\|^{2}\Bigr).
\end{aligned}
\tag{13}
\]

\noindent\textbf{[Step 7]} By $L$‑smoothness,
\[
\|\nabla f(x_k)\|^{2}\le 2L\bigl(f(x_k)-f(x^\star)\bigr).
\tag{14}
\]

\noindent\textbf{[Step 8]} Insert (14) into (13) and collect terms:
\[
\begin{aligned}
\E\!\left[\|x_{k+1}-x^\star\|^{2}\mid\mathcal{F}_k\right]
&\le\Bigl(1-\alpha\mu+2\alpha^{2}L^{2}\Bigr)\|x_k-x^\star\|^{2}
+\frac{2\alpha^{2}L^{2}}{n}\sum_{i=1}^{n}\|\phi_i^{k}-x^\star\|^{2}\\
&\quad -2\alpha\bigl(1-\alpha L\bigr)\bigl(f(x_k)-f(x^\star)\bigr).
\end{aligned}
\tag{15}
\]

\noindent\textbf{[Step 9]} Update the table entries (3). For the sampled index $i_k$ we have
\[
\|\phi_{i_k}^{k+1}-x^\star\|^{2}
= \|x_k-x^\star\|^{2},
\]
while for all $j\neq i_k$, $\phi_j^{k+1}=\phi_j^{k}$. Taking expectation over $i_k$ yields
\[
\E\!\left[\frac1n\sum_{j=1}^{n}\|\phi_j^{k+1}-x^\star\|^{2}\mid\mathcal{F}_k\right]
= \frac1n\|x_k-x^\star\|^{2}+\Bigl(1-\frac1n\Bigr)\frac1n\sum_{j=1}^{n}\|\phi_j^{k}-x^\star\|^{2}.
\tag{16}
\]

\noindent\textbf{[Step 10]} Form the Lyapunov function
\[
\Phi_k:=\|x_k-x^\star\|^{2}+\frac{\alpha}{nL}\sum_{i=1}^{n}\|\phi_i^{k}-x^\star\|^{2}.
\]
Combine (15) and (16) with the coefficient $\frac{\alpha}{nL}$ for the table part:
\[
\begin{aligned}
\E[\Phi_{k+1}\mid\mathcal{F}_k]
&\le\Bigl(1-\alpha\mu\Bigr)\|x_k-x^\star\|^{2}
+\frac{\alpha}{nL}\Bigl(1-\frac1n\Bigr)\sum_{i=1}^{n}\|\phi_i^{k}-x^\star\|^{2}\\
&\quad +\frac{2\alpha^{2}L^{2}}{n}\sum_{i=1}^{n}\|\phi_i^{k}-x^\star\|^{2}
-2\alpha\bigl(1-\alpha L\bigr)\bigl(f(x_k)-f(x^\star)\bigr).
\end{aligned}
\tag{17}
\]

\noindent\textbf{[Step 11]} Choose $\alpha\le\frac{1}{3L}$.
Then $1-\alpha L\ge\frac23$ and $2\alpha^{2}L^{2}\le\frac{2}{9}\alpha$.
Consequently,
\[
\frac{2\alpha^{2}L^{2}}{n}\sum_{i=1}^{n}\|\phi_i^{k}-x^\star\|^{2}
\le \frac{\alpha}{nL}\cdot\frac{2}{9}\sum_{i=1}^{n}\|\phi_i^{k}-x^\star\|^{2}.
\tag{18}
\]
Moreover,
\[
\Bigl(1-\frac1n\Bigr)+\frac{2}{9}\le 1,
\]
so the coefficient of the table term in (17) does not exceed $\frac{\alpha}{nL}$.

\noindent\textbf{[Step 12]} Dropping the non‑negative term $-2\alpha(1-\alpha L)(f(x_k)-f(x^\star))$ yields the clean contraction
\[
\E[\Phi_{k+1}\mid\mathcal{F}_k]\le\bigl(1-\alpha\mu\bigr)\Phi_k .
\tag{19}
\]

\noindent\textbf{[Step 13]} Taking full expectation and iterating (19) gives
\[
\E[\Phi_{k}]\le\bigl(1-\alpha\mu\bigr)^{k}\Phi_{0}.
\tag{20}
\]

\noindent\textbf{[Step 14]} Since $f$ is $L$‑smooth,
\[
f(x_k)-f(x^\star)\le\frac{L}{2}\|x_k-x^\star\|^{2}\le\frac{L}{2}\Phi_k .
\]
Combine with (20) to obtain (5), completing the proof. \hfill $\square$

\section*{Rate Derivation (With Constants)}
From Theorem \ref{thm:linear} and the steps above:
\[
\boxed{\;
\E\bigl[f(x_k)-f(x^\star)\bigr]
\le \frac{L}{2}\bigl(1-\alpha\mu\bigr)^{k}
\Bigl(\|x_0-x^\star\|^{2}+\frac{\alpha}{nL}\sum_{i=1}^{n}\|\phi_i^{0}-x^\star\|^{2}\Bigr)
\;}
\]
Choosing the maximal admissible stepsize $\alpha=1/(3L)$ gives the explicit contraction factor $1-\mu/(3L)$.

\section*{Special Cases}
\begin{itemize}
\item \textbf{Purely convex ($\mu=0$).} The Lyapunov contraction (19) reduces to a bound on the sum of squared distances. Applying Jensen’s inequality to the average iterate $\bar x_T$ yields the sublinear rate (6).
\item \textbf{Finite data with $n=1$.} SAGA collapses to classic gradient descent with stepsize $\alpha\le 1/L$, and (5) reproduces the standard linear convergence of GD.
\item \textbf{Large $n$ limit.} As $n\to\infty$, the table term becomes negligible; the algorithm behaves like SGD with a variance‑reduced gradient, preserving the $O(1/T)$ rate in the convex case.
\end{itemize}

\end{document}